% Options for packages loaded elsewhere
\PassOptionsToPackage{unicode}{hyperref}
\PassOptionsToPackage{hyphens}{url}
\documentclass[
  letterpaper,
]{article}
\usepackage{xcolor}
\usepackage[margin=0.8in]{geometry}
\usepackage{amsmath,amssymb}
\setcounter{secnumdepth}{-\maxdimen} % remove section numbering
\usepackage{iftex}
\ifPDFTeX
  \usepackage[T1]{fontenc}
  \usepackage[utf8]{inputenc}
  \usepackage{textcomp} % provide euro and other symbols
\else % if luatex or xetex
  \usepackage{unicode-math} % this also loads fontspec
  \defaultfontfeatures{Scale=MatchLowercase}
  \defaultfontfeatures[\rmfamily]{Ligatures=TeX,Scale=1}
\fi
\usepackage{lmodern}
\ifPDFTeX\else
  % xetex/luatex font selection
\fi
% Use upquote if available, for straight quotes in verbatim environments
\IfFileExists{upquote.sty}{\usepackage{upquote}}{}
\IfFileExists{microtype.sty}{% use microtype if available
  \usepackage[]{microtype}
  \UseMicrotypeSet[protrusion]{basicmath} % disable protrusion for tt fonts
}{}
\makeatletter
\@ifundefined{KOMAClassName}{% if non-KOMA class
  \IfFileExists{parskip.sty}{%
    \usepackage{parskip}
  }{% else
    \setlength{\parindent}{0pt}
    \setlength{\parskip}{6pt plus 2pt minus 1pt}}
}{% if KOMA class
  \KOMAoptions{parskip=half}}
\makeatother
\usepackage{listings}
\newcommand{\passthrough}[1]{#1}
\lstset{defaultdialect=[5.3]Lua}
\lstset{defaultdialect=[x86masm]Assembler}
\setlength{\emergencystretch}{3em} % prevent overfull lines
\providecommand{\tightlist}{%
  \setlength{\itemsep}{0pt}\setlength{\parskip}{0pt}}
% Preserve readable code, command, and table layout in Pandoc-generated review PDFs.
\usepackage{etoolbox}
\usepackage{pdflscape}
\lstset{
  basicstyle=\ttfamily\footnotesize,
  breaklines=true,
  breakatwhitespace=false,
  columns=fullflexible,
  keepspaces=true,
  showstringspaces=false,
  upquote=true
}
\AtBeginEnvironment{longtable}{\footnotesize}
\renewcommand{\arraystretch}{1.08}
\setlength{\emergencystretch}{3em}
\usepackage{bookmark}
\IfFileExists{xurl.sty}{\usepackage{xurl}}{} % add URL line breaks if available
\urlstyle{same}
% fallback for those not using the hyperref driver hyperxmp:
\makeatletter
\@ifundefined{xmpquote}{\newcommand{\xmpquote}[1]{#1}}{}
\makeatother
\hypersetup{
  pdftitle={Password Bonobo Python Reimplementation - Design Specification},
  hidelinks,
  pdfcreator={LaTeX via pandoc}}

\title{Password Bonobo Python Reimplementation - Design Specification}
\author{}
\date{}

\begin{document}
\maketitle

{
\setcounter{tocdepth}{3}
\tableofcontents
}
\section{Password Bonobo Python Reimplementation - Design
Specification}\label{password-bonobo-python-reimplementation---design-specification}

Date: 2026-08-22

Status: Approved design, pending implementation planning

Supersedes: The native Tcl implementation target in \href{./password-bonobo-url-audit-design.md}{Password Bonobo URL
Audit and Cleanup}. The approved URL-audit behavior and safety requirements remain authoritative except where this
specification replaces the platform architecture.

\subsection{1. Purpose}\label{1-purpose}

Password Bonobo will be a modern, cross-platform password manager with full meaningful Password Gorilla feature
compatibility and a defining URL-audit and credential-cleanup workflow.

The application will be implemented around a fully typed Python core with platform-appropriate clients. It will preserve
PasswordSafe V3 interoperability, remain local-file-first, and avoid a Bonobo-hosted cloud service.

The project will be an original, idiomatic implementation rather than a function-by-function Tcl translation.

\subsection{2. Approved Product Decisions}\label{2-approved-product-decisions}

\begin{itemize}
\tightlist
\item
  Preserve meaningful Password Gorilla features and database behavior without reproducing its visual layout.
\item
  Modernize the user experience on every platform.
\item
  Support Windows, macOS, Linux, and Android as tier-one platforms.
\item
  Support ChromeOS primarily through the Android application, with the Linux application as an optional fallback.
\item
  Treat iOS as an explicit supported target delivered after the secure core and Android path are proven.
\item
  Keep BSD-family portability in the core, with official support dependent on packaging, continuous integration, and
  real-system testing.
\item
  Use a fully typed, UI-independent Python core.
\item
  Permit narrowly scoped Kotlin and Swift code for native user interfaces and operating-system integrations.
\item
  Keep vault storage local-file-first and use operating-system document providers for user-selected synchronization
  locations.
\item
  Operate no Bonobo account system, cloud vault, or synchronization service.
\item
  Make loss of credentials or user-authored metadata unacceptable.
\item
  Treat derived audit results as replaceable state.
\item
  Use GPL-3.0-or-later as the intended project license, subject to dependency and distribution review.
\end{itemize}

\subsection{3. Licensing and Source Provenance}\label{3-licensing-and-source-provenance}

Password Gorilla declares GPL-2.0-or-later licensing. Bonobo may study a pinned Gorilla revision for compatibility
research, but no Gorilla implementation will enter a Bonobo product build.

The compatibility workflow will use these boundaries:

\begin{enumerate}
\def\labelenumi{\arabic{enumi}.}
\tightlist
\item
  Preserve the selected upstream Gorilla revision as an untouched, read-only research reference outside the Bonobo
  implementation tree.
\item
  Record observable features, workflows, edge cases, and compatibility expectations in neutral prose and test cases.
\item
  Implement Bonobo from the approved specifications, official PasswordSafe format documentation, neutral compatibility
  dossier, and independently designed architecture.
\item
  Do not copy Gorilla source, comments, identifiers, file organization, control flow, UI assets, translations, or other
  copyrightable expression.
\item
  Maintain provenance records for imported fixtures, algorithms, dependencies, and third-party assets.
\end{enumerate}

Bonobo-authored code intended for iOS distribution will include a narrowly drafted App Store distribution exception. Any
code not covered by that permission will be excluded from the iOS build. The final exception, dependency set, and
current Apple terms require review before an iOS release.

Future contributions must be accepted under terms that preserve the project\textquotesingle s GPL license and any
published distribution exception. The repository will document those terms before accepting external contributions.

\subsection{4. Compatibility Contract}\label{4-compatibility-contract}

Bonobo will treat a PasswordSafe V3 file as a lossless document rather than a simplified collection of known values.

For each header and record, the core will retain:

\begin{itemize}
\tightlist
\item
  Parsed standard fields.
\item
  Stable database and record UUIDs.
\item
  Field ordering and representation where compatibility requires preservation.
\item
  Every unknown field\textquotesingle s type identifier and original plaintext field bytes.
\item
  The distinction between standard data, user-authored metadata, and derived application state.
\end{itemize}

Editing a known value will change only the intended semantic field. Re-encryption will naturally change ciphertext, but
unknown field type and data bytes must survive the read-write cycle unchanged.

Bonobo will read supported older PasswordSafe V3 files without requiring a format migration. It will not silently raise
a file\textquotesingle s declared format level merely because Bonobo opened or saved it.

New user-authored metadata may use standardized custom fields with a Bonobo namespace only after actual Gorilla and
Password Safe round-trip tests prove that the claimed client versions preserve those fields. When a feature cannot meet
that requirement, Bonobo will store it in an optional separately encrypted metadata document or disable it in strict
interoperability mode. Important user data will never be placed where a supported client can silently erase it.

Derived URL-audit state will remain in memory for the initial release. Persistent audit history is not required by this
design.

\subsection{5. System Architecture}\label{5-system-architecture}

\subsubsection{5.1 Python core}\label{51-python-core}

The bonobo\_core package will own:

\begin{itemize}
\tightlist
\item
  PasswordSafe V3 parsing, serialization, authentication, and validation.
\item
  Vault sessions, records, groups, policies, history, aliases, shortcuts, and protected entries.
\item
  Unknown-field preservation and extension policy.
\item
  Application use cases and transactional operations.
\item
  URL collection, sanitization, scan orchestration, classification, and cleanup rules.
\item
  Archive creation and validation.
\item
  Provider-independent conflict detection and merge models.
\item
  Typed failure categories and safe diagnostic metadata.
\end{itemize}

The core will not import desktop, Android, iOS, or cloud-provider user-interface code.

\subsubsection{5.2 Platform clients}\label{52-platform-clients}

\begin{itemize}
\tightlist
\item
  Desktop will use PySide6 and Qt Quick on Windows, macOS, and Linux.
\item
  Android and ChromeOS will use Kotlin and Jetpack Compose.
\item
  iOS will use Swift and SwiftUI.
\end{itemize}

Android and iOS will embed CPython and call a narrow in-process application facade. The facade will expose task-oriented
operations and typed data-transfer objects. Secret values will cross the boundary only for explicit reveal, copy, edit,
Autofill, AutoType, or credential-provider actions.

\subsubsection{5.3 Platform services}\label{53-platform-services}

The core will define protocols for:

\begin{itemize}
\tightlist
\item
  Document selection, coordinated access, and provider publication.
\item
  Device-bound secure key storage and biometric authorization.
\item
  Clipboard handling and automatic clearing.
\item
  Browser launching.
\item
  HTTP transport.
\item
  Application lifecycle and idle locking.
\item
  File logging.
\item
  Accessibility-relevant notifications where needed.
\end{itemize}

Each native client will provide implementations without duplicating domain rules.

\subsubsection{5.4 Compatibility assets}\label{54-compatibility-assets}

The compatibility area will contain:

\begin{itemize}
\tightlist
\item
  The neutral Gorilla behavior dossier.
\item
  A feature-parity matrix.
\item
  Synthetic PasswordSafe fixtures.
\item
  Cross-client round-trip expectations.
\item
  Reproducible compatibility test instructions.
\end{itemize}

No production or test fixture may contain real credentials or personal vault data.

\subsection{6. Document Lifecycle}\label{6-document-lifecycle}

\subsubsection{6.1 Opening}\label{61-opening}

\begin{enumerate}
\def\labelenumi{\arabic{enumi}.}
\tightlist
\item
  The native client asks the operating system to select a PasswordSafe document.
\item
  The provider adapter obtains a complete encrypted snapshot and any available provider revision metadata.
\item
  Bonobo calculates a SHA-256 fingerprint of the encrypted bytes.
\item
  Mobile clients place the encrypted snapshot in an app-private working location for offline access.
\item
  The core authenticates the master password, validates the file, and opens an in-memory session.
\item
  The client receives non-secret summaries until the user explicitly requests a sensitive operation.
\end{enumerate}

\subsubsection{6.2 Ordinary edits}\label{62-ordinary-edits}

\begin{enumerate}
\def\labelenumi{\arabic{enumi}.}
\tightlist
\item
  The user confirms a complete edit form.
\item
  The core validates the proposed record and updates the in-memory model.
\item
  Bonobo transactionally writes and validates the encrypted working copy.
\item
  The provider adapter verifies that the external revision still matches the opening baseline.
\item
  An unchanged external document may receive the validated new version.
\item
  An unavailable provider produces Sync pending without losing the encrypted local edit.
\item
  A changed provider document produces Conflict and is never overwritten automatically.
\end{enumerate}

\subsubsection{6.3 URL-audit cleanup}\label{63-url-audit-cleanup}

Audit deletions remain a staged transaction even though ordinary edits commit when their edit form is confirmed.
Affected rows remain visibly marked Deleted (unsaved) or Archived \& deleted (unsaved) until the user explicitly saves.

Closing a session with staged cleanup requires the user to save or discard that cleanup. Discarding staged cleanup does
not roll back unrelated ordinary edits already committed by the user.

\subsubsection{6.4 Lock and close}\label{64-lock-and-close}

Locking, closing, replacing a vault, or losing authentication will:

\begin{itemize}
\tightlist
\item
  Cancel outstanding URL-audit requests.
\item
  Close or clear sensitive views.
\item
  Clear plaintext and key material as far as the runtime permits.
\item
  Retain only encrypted files and bounded non-sensitive state.
\end{itemize}

\subsection{7. Provider-Neutral File Synchronization}\label{7-provider-neutral-file-synchronization}

Bonobo will not implement provider-specific Google Drive, Dropbox, OneDrive, or iCloud accounts. Users will choose files
through native document-provider interfaces or synchronized desktop folders.

For each open document, Bonobo will remember outside the PasswordSafe file:

\begin{itemize}
\tightlist
\item
  The standard database UUID.
\item
  The SHA-256 fingerprint of the encrypted opening snapshot.
\item
  Provider revision, generation, or ETag information when available.
\item
  File size and modification time as secondary signals.
\end{itemize}

Bonobo will not add a custom PasswordSafe field solely for synchronization numbering. A custom counter cannot prevent
two offline writers from selecting the same next value and may be discarded by another client.

Before publication, Bonobo will compare the provider document with the recorded baseline. If it changed, Bonobo will
preserve both valid versions and offer explicit reload, save-as-copy, or UUID-based conflict resolution.

Provider rules include:

\begin{itemize}
\tightlist
\item
  Never rely on a file lock to coordinate different devices.
\item
  Never truncate a provider document in place.
\item
  Validate a complete new vault before provider replacement or upload.
\item
  Use native coordinated-access facilities when the platform provides them.
\item
  Keep bounded, encrypted recovery revisions.
\item
  Display Local changes, Sync pending, Conflict, and Provider unavailable honestly.
\item
  Never silently select a winner based only on timestamps.
\item
  Never place derived audit cache or diagnostic logs beside the synchronized vault.
\end{itemize}

Providers can observe filenames, sizes, and modification times. Product documentation will explain that metadata leakage
even though PasswordSafe encrypts vault contents.

\subsection{8. URL Audit and Cleanup}\label{8-url-audit-and-cleanup}

The behavior, classifications, review workflow, archive-first deletion transaction, and privacy requirements in the
existing URL-audit specification remain in force.

The Python architecture divides responsibility as follows:

\begin{itemize}
\tightlist
\item
  The core sanitizes URLs, strips query strings and fragments, detects hygiene issues, schedules bounded work, evaluates
  redirects, performs root fallback decisions, detects parking indicators, and classifies observations.
\item
  Desktop uses an asynchronous Python HTTP transport.
\item
  Android uses the native Android HTTP and TLS stack through Kotlin.
\item
  iOS uses URLSession through Swift.
\end{itemize}

Every transport will return the same normalized observation model, including status, redirect chain, bounded response
prefix, and categorized DNS, connection, timeout, or TLS failure.

Only the sanitized audit URL may cross the transport boundary. Credentials, record identifiers, group names, titles,
database identity, query values, fragments, and other sensitive record data must not enter requests or HTTP headers.

Results appear incrementally. Cancellation stops new work and cancels outstanding requests where practical. Locking or
replacing the vault cancels the scan and destroys sensitive result state.

The original stored URL remains separate from the audit URL and is used only for an explicit browser-opening action.
Archive and deletion actions call core transactions; user-interface code cannot directly delete records or construct
archive files.

\subsection{9. Security Model}\label{9-security-model}

The threat model includes:

\begin{itemize}
\tightlist
\item
  A stolen locked device.
\item
  An untrusted file-hosting provider.
\item
  Malicious or compromised websites.
\item
  Stale synchronized files and concurrent writers.
\item
  Sensitive clipboard contents.
\item
  Diagnostic leakage.
\item
  Accidental destructive actions.
\end{itemize}

Bonobo will:

\begin{itemize}
\tightlist
\item
  Use reviewed cryptographic implementations and documented PasswordSafe algorithms.
\item
  Never invent encryption primitives.
\item
  Never store the master password in plaintext.
\item
  Keep decrypted vault data inside an authenticated session.
\item
  Minimize secret copies and clear mutable buffers where practical.
\item
  Document that CPython cannot guarantee perfect memory zeroization.
\item
  Disable telemetry by default.
\item
  Avoid executing stored URLs, commands, or data without an explicit supported user action.
\end{itemize}

\subsubsection{9.1 Biometric unlock}\label{91-biometric-unlock}

Biometric unlock is optional and explicitly enabled per vault and device:

\begin{enumerate}
\def\labelenumi{\arabic{enumi}.}
\tightlist
\item
  The first unlock requires the master password.
\item
  Bonobo prepares the minimum vault-unlock material required for later access.
\item
  A device-bound hardware-backed key wraps that material through Android Keystore or Apple security facilities.
\item
  Successful biometric or device-owner authentication is required to unwrap it.
\item
  Wrapped material never enters the PasswordSafe document or synchronization provider.
\item
  Master-password changes and relevant device-security changes invalidate or replace it.
\item
  Users can revoke biometric unlock and require the master password again.
\end{enumerate}

Biometric unlock may operate across application restarts after successful enrollment.

\subsubsection{9.2 Autofill and credential providers}\label{92-autofill-and-credential-providers}

Android Autofill and the iOS credential-provider extension will be native components. They may receive only the minimum
matching metadata and explicitly requested credential values from an authenticated core session.

These components cannot enumerate decrypted vault contents independently, initiate URL audits, bypass lock policy, or
publish a vault without the same validation and conflict controls as the main application.

\subsubsection{9.3 Clipboard and diagnostics}\label{93-clipboard-and-diagnostics}

Copied secrets will use platform-sensitive clipboard markers and configurable automatic clearing where available.

Runtime logs will be written under the application sandbox\textquotesingle s logs directory. Formatted records will use
only INFO, DEBG, WARN, and CRIT severities and will include a timestamp, source, and safe message.

Logs must not contain credentials, tokens, full URLs, URL paths, query values, vault filenames, entry names, UUIDs, or
other sensitive database content.

\subsection{10. Failure Handling}\label{10-failure-handling}

The core will expose typed failure categories instead of platform exception text.

\begin{itemize}
\tightlist
\item
  Authentication, parsing, validation, archive, and provider failures leave the source vault unchanged.
\item
  Interrupted writes preserve the last validated vault and an encrypted recovery copy.
\item
  External changes block publication and enter conflict resolution.
\item
  Audit failures change classifications only and never mutate records.
\item
  Unsupported or malformed content fails closed while retaining the original encrypted bytes.
\item
  Unexpected termination may leave an encrypted pending working copy but never plaintext recovery data.
\item
  User-facing errors explain the safe next action without exposing sensitive diagnostics.
\end{itemize}

\subsection{11. Verification Strategy}\label{11-verification-strategy}

Core verification will include:

\begin{itemize}
\tightlist
\item
  Strict static type checking.
\item
  Project formatter and linter enforcement.
\item
  Unit tests for domain rules and application use cases.
\item
  Property-based tests for parsing, serialization, and transformations.
\item
  Fuzz testing of PasswordSafe parsing and malformed fields.
\item
  Golden synthetic vaults covering all supported standard and unknown fields.
\item
  Byte-preservation and semantic round trips through Bonobo, Gorilla, and Password Safe.
\item
  Archive reopen and identity verification before deletion.
\item
  Large-vault performance and memory tests.
\end{itemize}

URL-audit verification will include deterministic local HTTP services and shared transport contract tests for:

\begin{itemize}
\tightlist
\item
  Successful responses and redirects.
\item
  Cross-host redirects.
\item
  Authentication and rate-limit responses.
\item
  Live roots with dead saved paths.
\item
  DNS, connection, timeout, and TLS failures.
\item
  Redirect loops and limits.
\item
  Parked and ambiguous content.
\item
  Cancellation and bounded concurrency.
\end{itemize}

Provider verification will inject:

\begin{itemize}
\tightlist
\item
  Offline operation.
\item
  Interrupted upload and replacement.
\item
  Provider permission revocation.
\item
  External mutation immediately before save.
\item
  Conflicted copies and version conflicts.
\item
  Disk-full and quota failures.
\item
  Non-atomic provider behavior.
\end{itemize}

Platform verification will include:

\begin{itemize}
\tightlist
\item
  Continuous integration on Windows, macOS, and Linux.
\item
  Android emulator and physical-device tests.
\item
  Android Keystore, biometric, Autofill, background termination, and ChromeOS tests.
\item
  iOS simulator and physical-device tests.
\item
  Apple document coordination, Keychain, biometric, and credential-provider tests.
\item
  Accessibility, keyboard navigation, screen-reader, and responsive-layout tests.
\end{itemize}

Every release requires formatter, type checker, static security checks, dependency and license audit, packaging
validation, and the applicable automated and physical-device suites.

\subsection{12. Delivery Decomposition}\label{12-delivery-decomposition}

The program will use separate specifications and implementation plans for these subprojects:

\begin{enumerate}
\def\labelenumi{\arabic{enumi}.}
\tightlist
\item
  Repository foundation and compatibility dossier.
\item
  Lossless PasswordSafe core.
\item
  Vault application core and desktop foundation.
\item
  URL audit and cleanup.
\item
  Provider-safe files and conflict resolution.
\item
  Android and ChromeOS.
\item
  iOS.
\item
  Parity closure and stable release.
\end{enumerate}

\subsubsection{12.1 Repository foundation and compatibility
dossier}\label{121-repository-foundation-and-compatibility-dossier}

\begin{itemize}
\tightlist
\item
  Establish Git, repository policy, GPL licensing, App Store exception planning, contribution terms, CI skeleton, and
  durable project memory.
\item
  Pin the upstream Gorilla reference without placing it in the product tree.
\item
  Produce the neutral behavior dossier and feature-parity matrix.
\end{itemize}

\subsubsection{12.2 Lossless PasswordSafe core}\label{122-lossless-passwordsafe-core}

\begin{itemize}
\tightlist
\item
  Implement typed parsing, writing, validation, unknown-field preservation, and transactional local files.
\item
  Establish conformance, fuzz, and round-trip suites before production UI work depends on the core.
\end{itemize}

\subsubsection{12.3 Vault application core and desktop
foundation}\label{123-vault-application-core-and-desktop-foundation}

\begin{itemize}
\tightlist
\item
  Implement sessions, records, search, groups, policies, history, protected entries, archive operations, and platform
  protocols.
\item
  Deliver the modern desktop shell and packaging foundation.
\end{itemize}

\subsubsection{12.4 URL audit and cleanup}\label{124-url-audit-and-cleanup}

\begin{itemize}
\tightlist
\item
  Adapt the approved audit design to Python classification and native mobile transports.
\item
  Deliver review, archive, staged deletion, cancellation, and safety testing.
\end{itemize}

\subsubsection{12.5 Provider-safe files and conflict resolution}\label{125-provider-safe-files-and-conflict-resolution}

\begin{itemize}
\tightlist
\item
  Implement native document providers, encrypted working copies, pending publication, conflict detection, recovery, and
  explicit merge workflows.
\end{itemize}

\subsubsection{12.6 Android and ChromeOS}\label{126-android-and-chromeos}

\begin{itemize}
\tightlist
\item
  Deliver the Compose client, embedded Python core, provider integration, biometrics, Autofill, lifecycle safety,
  offline operation, and ChromeOS qualification.
\item
  Android is not complete until biometric unlock and Autofill pass physical-device tests.
\end{itemize}

\subsubsection{12.7 iOS}\label{127-ios}

\begin{itemize}
\tightlist
\item
  Deliver the SwiftUI client, embedded Python core, coordinated documents, Keychain and biometrics, credential provider,
  licensing audit, and App Store preparation.
\end{itemize}

\subsubsection{12.8 Parity closure and stable release}\label{128-parity-closure-and-stable-release}

\begin{itemize}
\tightlist
\item
  Close remaining Gorilla compatibility gaps.
\item
  Complete accessibility, performance, migration documentation, and release audits.
\end{itemize}

\subsection{13. Program Acceptance Criteria}\label{13-program-acceptance-criteria}

The program-level design is satisfied when:

\begin{enumerate}
\def\labelenumi{\arabic{enumi}.}
\tightlist
\item
  Existing supported Gorilla PasswordSafe V3 databases open without losing standard, unknown, or user-authored fields.
\item
  Bonobo saves pass the documented cross-client round-trip matrix.
\item
  Windows, macOS, and Linux provide a modern interface with meaningful Gorilla feature compatibility.
\item
  URL audit and cleanup meet the approved classification, privacy, archive, and deletion-safety criteria.
\item
  External document providers cannot cause silent last-writer-wins overwrites.
\item
  Ordinary edits survive mobile lifecycle termination in encrypted form.
\item
  Bulk audit cleanup remains staged until explicit save.
\item
  Android supports secure offline use, biometric unlock, Autofill, and qualified ChromeOS operation.
\item
  iOS uses only code and dependencies eligible for its approved distribution terms.
\item
  No Bonobo cloud service or account is required.
\item
  Security-sensitive logs, network requests, caches, fixtures, and recovery files contain no prohibited plaintext data.
\item
  Every claimed platform passes its automated, packaging, and required physical-device release gates.
\end{enumerate}

\subsection{14. Known Risks}\label{14-known-risks}

\begin{itemize}
\tightlist
\item
  Cross-client unknown-field preservation may differ from the PasswordSafe specification and must be verified
  empirically.
\item
  Python embedding and binary dependency packaging are more complex on Android and iOS than on desktop.
\item
  CPython cannot guarantee complete in-memory secret zeroization.
\item
  App Store terms and required exceptions can change before the iOS release.
\item
  Third-party document providers expose inconsistent coordination and conflict capabilities.
\item
  Multiple native presentation layers increase implementation and testing cost.
\item
  Full Gorilla compatibility is a multi-release effort and requires a continuously maintained parity matrix.
\end{itemize}

These risks justify the phased delivery model and do not relax data-integrity or security requirements.

\subsection{15. References}\label{15-references}

\begin{itemize}
\tightlist
\item
  Password Gorilla: \url{https://github.com/zdia/gorilla}
\item
  PasswordSafe V3 format: \url{https://github.com/pwsafe/pwsafe/blob/master/docs/formatV3.txt}
\item
  CPython on Android: \url{https://docs.python.org/3.14/using/android.html}
\item
  CPython on iOS: \url{https://docs.python.org/3.14/using/ios.html}
\item
  Qt supported platforms: \url{https://doc.qt.io/qtforpython-6/overviews/qtdoc-supported-platforms.html}
\item
  Android Storage Access Framework: \url{https://developer.android.com/guide/topics/providers/document-provider}
\item
  Apple shared document coordination: \url{https://developer.apple.com/documentation/technologyoverviews/shared-data}
\item
  GNU GPLv3 guide: \url{https://www.gnu.org/licenses/quick-guide-gplv3.html}
\item
  U.S. Copyright Office Circular 61: \url{https://www.copyright.gov/circs/circ61.pdf}
\end{itemize}

\end{document}
